\documentclass[10pt]{article}

\usepackage[utf8]{inputenc}

\usepackage[T1]{fontenc}

\usepackage{amsmath}

\usepackage{amsfonts}

\usepackage{amssymb}

\usepackage[version=4]{mhchem}

\usepackage{stmaryrd}

\usepackage{bbold}



\DeclareUnicodeCharacter{0131}{$\imath$}



\begin{document}

распределению (5) при больших $V$, близко к среднему $\langle f\rangle \bar{\xi}^{0}, \ldots . \bar{\xi} k$, вычисленному по микроканонич. ансамблю на поверхности $S \bar{\xi}^{0}, \ldots, \bar{\xi}^{k}$, где $\left\langle\bar{\xi}^{i}\right\rangle=\left\langle H_{V}^{i}\right\rangle$ $>\beta, \mu_{1}, \ldots, \mu_{k}$. Доказательство этой әквивалентности также составляет одну из общих математич. проблем статистич. механики и термодинамики (см. [5], [6], [7]).



Принятый в статистич. механике способ описания систем оправдан при достаточно большом объеме области $V$, иначе говоря, статистич. механика изучает асимптотич, свойства систем в предельном переходе $\mathrm{V} \nearrow \mathbb{R}^{3}$ (т. е. рассматривается нек-рая последователь ность систем из одних и тех же частиц, заключенных соответственно в объемах $V_{1} \subset V_{2} \subset$, . ., причем $U_{n} V_{n}=$ $=\mathbb{R}^{3}$ ). Этот предельный переход наз. обычно т е р м один а мит. п р е де льным пе р ех о дом. Одна из первых задач, связанных с термодинамич. пределом, состоит в том, чтобы, исходя из равновесных ансамблей, определить т. н. те р м о д и н а м и ч. пот енциалы и соотношения термод и н а м и к и. Оказывается, все термодинамич. потенциалы могут быть найдены из асимптотики при $V \nearrow \mathbb{R}^{3}$ нормировочных множителей $Q^{-1}, Z^{-1}, \tilde{Z}^{-1}$ и т. Д. в ансамблях (3), (5), (6), так, напр., т е р м один а мич. по тен ци а л Ги б б с а равен



\[

p\left(\beta, \mu_{1}, \ldots, \mu_{k}\right)=\lim _{V \nearrow \mathbb{R}^{3}} \frac{\ln z}{|V|},

\]



где $Z^{-1}$ - нормирующий множитель в гиббсовском ансамбле (5). Аналогично вводятся и др. термодинамич. функции, а также устанавливаются связывающие их соотношения. Большинство возникающих здесь математич. задач (существование предела, свойства термодинамич. потенциалов и т. д.) исследовано довольно полно, хотя и имеется ряд нерешенных вопросов (см., напр., [7]).



С конца 60-х гг. 20 в. в математич. статистич. механике утвердился следующий общий подход: вместо изучения асимптотич. свойств конечных систем в термодинамическом предельном переходе следует рассматривать определенным образом построенные идеализированные бесконечные системы, характеристики к-рых совпадают с исследуемой асимптотикой (неявно такая точка зрения встречалась и в более ранних работах). Рассмотрение бесконечных систем придает наглядный смысл несколько формальной процедуре термодинамического предельного перехода и позволяет вообще обойтись без нее. Фазовое пространство $\Omega_{\infty}$ бесконечной системы состоит из бесконечных конфигураций частиц $\omega=\left\{x_{1}, x_{2}, \ldots\right\}, x_{i} \in X^{V}, i=1,2, \ldots$, располагающихся во всем пространстве $\mathbb{R}^{3}$, a их дйнамика $U_{t}^{\infty}: \Omega_{\infty} \rightarrow \Omega_{\infty}, t \in \mathbb{R}$, строится как предел динамик $U_{t}^{V}$ конечных систем при $V \nearrow \mathbb{R}^{3}$. Макроскопич. состояния бесконечной системы задаются по-прежнему вероятностными распределениями на пространстве $\Omega_{\infty}$, к-рые эволюционируют в соответствии с динамикой $U_{t}^{\infty}$ в $\Omega_{\infty}$ (см. (1)). На пространстве $\Omega_{\infty}$ можно ввести предельные гиббсовские распределения $p_{\beta, \mu_{1}}^{\infty} \ldots, \mu_{k}$, к-рые строятся определенным образом с помощью гиббсовских распределений (5) $p_{\beta, \mu_{1}}^{V}, \ldots, \mu_{k}$ в конечных системах (см. [5], [9]). Хотя введение бесконечных систем является общепринятым и плодотворным ириемом, оно приводит к сложным и во многом нерешенным (1984) еще математич. задачам. Сложным, напр., оказывается построение динамики $U_{t}^{\infty}$, построение предельных гиббсовских распределений, исследование их свойств і т. Д.



Одна из главных ироблем статистич. механики состоит в изучении т. н. ф а з о в т. е. резкого изменения состояния макроскопич. системы, находящейся в состоянии равновесия, при небольшом изменении описывающих это равновесие параметров - температуры, плотности частиц, давления и т. д. При современном математич. подходе в терминах предельных гиббсовских распределений задачу о фазовых переходах можно описать следующим образом: при нек-рых значениях параметров $\beta, \mu_{1}$, ..., $\mu_{k}$ можно построить, вообще говоря, несколько гиббсовских распределений на $\Omega_{\infty}$, инвариантных относительно действия группы $T^{3}$ сдвигов в пространстве $\mathbb{R}^{3}$ (или нек-рой ее подгруппы $G \subset T^{3}$ такой, что факторгруппа $T^{3} / G$ компактна), и эргодичных относительно әтой группы (т. н. ч и с т ы е фа зы). Точка $\left(\beta, \mu_{1}, \ldots, \mu_{k}\right)$ пространства параметров наз. регулярной, если существует достаточно малая ее окрестность, внутри к-рой структура множества чистых фаз, а также основные их качественные свойства (напр., характер убывания корреляций) остаются неизменными. При әтом предполагается, что все числовые характеристики этих распределений (корреляционные функции, семиинварианты и т. д.) в окрестности регулярных точек зависят от параметров $\beta, \mu_{1}, \ldots$ ..., $\mu_{k}$ аналитически. Все остальные (не регулярные) точки в пространстве параметров $\beta, \mu_{1}, \ldots ., \mu_{k}$ и являются точками фазового перехода. Таким образом, в точках фазового перехода происходит резкое изменение либо в структуре гиббсовских распределений (скажем, исчезает или возникает новая фаза) или в их свойствах (напр., убывание корреляций из экспоненциального становится степенным). При этом считается, что какие-нибудь из характеристик распределения как функции параметров $\beta, \mu_{1}, \ldots, \mu_{k}$ имею' в точке фазового перехода особенность. Описать для каждой конкретной системы структуру фаз, их свойства, определить точки фазового перехода, характер особенностей в этих точках и т. Д.- этот круг вопросов и составляет проблему фазовых переходов. Хотя существует большой класс модельных систем, для к-рых (при малых значениях температуры) разработаны нек-рые общие методы решения этой задачи (см. [9]), теория фазовых переходов еще далека от окончательного завершения. Особенно сложным является изучение т. н. критич. точек (в к-рых, грубо говоря, шроисходит слияние разных фаз; см. [10]), поскольку в этих точках гиббсовское распределение имеет очень медленное убывание корреляций.



Большой круг проблем статистич. механики связан с изучением эволюций распределений на фазовом пространстве, в частности с проблемой р е л а к с а ц и и, т. е. приближения к равновесию. Считается, что по прошествии большого времени всякое распределение на фазовом пространстве приближается к равновесному (гиббсовскому) распределению. Несмотря на то, что выработано много общих представлений о механизме этого процесса, а также исследован ряд упрощенных его моделей, законченной теории пока (1984) не существует. Основные представления о процессе релаксации (во многом еще гипотетичные) сводятся к тому, что этот процесс проходит три стадии. На первой из них (за время столкновений нескольких частиц) расıределение $p_{t}$ приходит к такому режиму әволюции, к-рый целиком определяется изменением первой корреляционной функции (т. е. распределением в одночастичном пространстве $X)$. Затем, во второй кинетич. стадии, протекаюей за промежуток времени порядка продолжительности «свободного пробега" частицы, изменение первой корреляционной функции переходит в такой режим әволюции, при к-ром все зависит липь от средних значений плотности частиц, их скорости, плотности, энергии и т. д. Наконец, настунает последняя (гидродинамическая) стадия, во





\end{document}